\tikzstyle{confnode}=[draw, rounded corners=2pt, thick, fill=white,
                      inner sep=2.5pt, font=\small]
\tikzstyle{confcurrent}=[confnode, fill=red!20]
\tikzstyle{confedge}=[thick]
\tikzstyle{resultbox}=[draw, rounded corners=3pt, thick,
                       fill=green!25, inner sep=5pt, font=\normalsize]

% -------------------------------------------------
% Vstupní slovo pro NFA
% -------------------------------------------------
\newcommand{\NFAInputWide}[3]{%
    \node[font=\large] at (3.2,1.9) {$#1\textcolor{red}{#2}#3$};
}

\newcommand{\NFAInputWideFinished}[1]{%
    \node[font=\large] at (3.2,1.9) {$#1\textcolor{red}{\lambda}$};
}

% -------------------------------------------------
% Jeden slide animace pro NFA rozpoznávající podslovo 101
% #1 ... styl q0
% #2 ... styl q1
% #3 ... styl q2
% #4 ... styl q3
% #5 ... kód pro vstupní slovo
% #6 ... kód stromu výpočtu (+ případně výsledek)
% -------------------------------------------------
\newcommand{\SubstringNFAFrame}[6]{%
\begin{frame}{Výpočet nedeterministického automatu}{Příklad}
\begin{tikzpicture}[>=stealth]

    % vstup
    #5

    % automat
    \initialstate[#1]{q0}{(0,0)}{$q_0$}{(-1.2,0)}{west}
    \state[#2]{q1}{(2.1,0)}{$q_1$}
    \state[#3]{q2}{(4.2,0)}{$q_2$}
    \acceptingstate[#4]{q3}{(6.3,0)}{$q_3$}

    % přechody
    \transitionarrow[loop above]{q0}{above}{$0,1$}{q0}
    \transitionarrow{q0}{above}{$1$}{q1}
    \transitionarrow{q1}{above}{$0$}{q2}
    \transitionarrow{q2}{above}{$1$}{q3}
    \transitionarrow[loop above]{q3}{above}{$0,1$}{q3}

    % strom výpočtu
    \node[anchor=west,font=\normalsize] at (-0.3,-1.45) {Strom výpočtu:};
    #6

\end{tikzpicture}
\end{frame}
}